% ==============================================================================
% PLANTILLA MAESTRA - FILOSOFIA DEL DERECHO / UnADM
% ==============================================================================
% Uso:
% 1. Duplica este archivo y renombra la copia por semana/actividad.
% 2. Actualiza documenttitle, documentsubtitle y documentsubject.
% 3. Lee la planeacion de la semana y NO pegues las instrucciones completas en el
%    producto final. Usa las instrucciones solo como lista de cumplimiento.
% 4. Identifica primero el producto solicitado: mapa, cuadro, glosario, tabla,
%    estudio de caso, reflexion, video, etc. El formato pedido manda sobre el
%    gusto personal por escribir ensayo.
% 5. Revisa bibliografia, toma notas, define categorias y redacta con postura.
% 6. Entrega en PDF con la nomenclatura indicada por la planeacion.
%
% Formato institucional aplicado:
% - Fuente tipo Helvetica.
% - Interlineado 1.5.
% - Citas autor-anio con natbib: usar \citep{clave} o \citet{clave}.
% - Referencias en bibliografia BibTeX.
% - Redaccion academica, clara, con analisis propio.
% - Productos visuales elaborados preferentemente en LaTeX/TikZ: tablas,
%   cuadros comparativos, mapas conceptuales/mentales, lineas de tiempo,
%   diagramas de flujo, esquemas, cronologias, matrices y organizadores.
%
% Criterios derivados de retroalimentacion docente:
% - No limitarse a reproducir fuentes: interpretar, conectar y tomar postura.
% - La pregunta guia antes de entregar debe ser: cual es mi posicion frente a lo
%   que expongo y que razones sostienen esa posicion.
% - Evitar redaccion excesivamente general, repeticion de estructuras y tono de
%   resumen automatico. Cada parrafo debe tener una decision propia.
% - Las fuentes no se mencionan como lista decorativa: cada autor debe sostener
%   una definicion, un criterio, una diferencia, una tension o una aplicacion.
% - Toda referencia final debe estar citada dentro del texto y toda cita dentro
%   del texto debe aparecer en referencias. Revisar coincidencia uno a uno.
% - Incluir conclusion personal cuando la actividad lo solicite.
% - En cuadros comparativos: categorias claras, jerarquias, secuencia y sintesis.
% - Revisar diseno solicitado: producto visual/cuadro, conclusiones y referencias.
% - Si la consigna pide mapa conceptual, cuadro, tabla o glosario, construir ese
%   producto. No convertirlo en texto expositivo salvo que la planeacion lo pida.
% - Usar APA 7 edicion en citas y referencias.
% - No presentar las indicaciones de la actividad dentro del documento final.
% - Si se usa IA como apoyo, el resultado debe estar leido, corregido, apropiado
%   y transformado con criterio personal; nunca entregar texto sin procesamiento.
%
% Protocolo contra observaciones de "IA" o "falta de autenticidad":
% - Abrir cada seccion con una idea propia, no con definiciones genericas.
% - Integrar por lo menos una frase de valoracion: "El punto critico es...",
%   "Desde esta perspectiva...", "La consecuencia juridica es...".
% - Pasar de definicion a funcion: que hace el concepto dentro del sistema.
% - Pasar de funcion a consecuencia: que ocurre si se interpreta mal.
% - En conclusiones, no repetir; responder que se comprendio, que se sostiene y
%   que cambia para la interpretacion o practica juridica.
%
% APA 7 operativo para esta materia:
% - La cita puede ser parafraseada o textual. Si la planeacion dice "explica con
%   tus propias palabras", preferir parafrasis con cita autor-anio.
% - En glosarios, colocar la cita junto a la definicion base y despues desarrollar
%   comentario propio. Formula: definicion base + cita -> funcion -> ejemplo o
%   consecuencia.
% - Usar cita textual solo cuando la frase exacta sea necesaria; incluir pagina:
%   \citep[p.~00]{clave}. No llenar todo el trabajo de comillas.
% - En citas de varias fuentes, no acumular autores sin distinguir aportes. Mejor
%   explicar que aporta cada fuente en una frase breve.
% - No usar URLs sueltas dentro del texto. Registrar la fuente en .bib y citarla.
%
% Criterios de redaccion personal (enfoque_personal.txt + guia estrategica):
% - No escribir para llenar espacio; escribir para posicionar una idea con peso.
% - Cada seccion debe responder: que es, como funciona, por que importa,
%   que implica y cual es la postura.
% - Evitar texto generico. Convertir informacion en criterio.
% - Estructura mental recomendada: problema -> sistema -> consecuencia -> postura.
% - Usar un tono formal, directo, analitico y con intencion.
% - Integrar pensamiento de sistemas: funcion, estructura, control, criterio,
%   costo, decision, consecuencia, limites, aplicacion real.
% - Usar primera persona solo con moderacion: "Desde esta perspectiva",
%   "Considero que", "La posicion que sostengo es". Evitar "yo creo".
%
% Checklist antes de entregar:
% [ ] El producto coincide exactamente con la tecnica solicitada en planeacion.
% [ ] El objetivo esta alineado con la planeacion.
% [ ] Las instrucciones no fueron copiadas como contenido del producto.
% [ ] Hay citas en APA autor-anio dentro del texto, ubicadas donde se usa la idea.
% [ ] La bibliografia final coincide con las citas del cuerpo, sin sobrantes.
% [ ] La introduccion abre con una afirmacion fuerte, no con relleno.
% [ ] El desarrollo explica, interpreta y conecta.
% [ ] Cada fuente cumple una funcion analitica, no solo aparece mencionada.
% [ ] Hay ejemplos juridicos o aplicacion al sistema mexicano cuando proceda.
% [ ] La conclusion incluye postura personal argumentada: posicion, razon y efecto.
% [ ] No aparecen instrucciones de la planeacion copiadas en el cuerpo final.
% [ ] Si la actividad pide producto visual, se construyo con TikZ/LaTeX o se
%     justifico por que se uso una captura externa.
% [ ] Si es mapa conceptual, hay jerarquia, palabras de enlace y relaciones.
% [ ] Si es cuadro comparativo, hay categorias comparables y sintesis critica.
% [ ] Si es glosario, cada concepto tiene definicion propia, cita y aplicacion.
% [ ] Si es caso, se identifican hechos, problema, conceptos, aplicacion y postura.
% [ ] Ortografia, acentos, sintaxis y nombres propios revisados.
% [ ] Si se uso IA como apoyo, el texto fue comprendido, revisado y apropiado.
% ==============================================================================

% CREACION DEL DOCUMENTO
\documentclass[
	spanish,
	letterpaper, oneside
]{article}

% ==============================================================================
% INFORMACION DEL DOCUMENTO
% Actualizar en cada actividad.
% Ejemplos:
% \def\documenttitle {Mapa conceptual de la Filosofia del Derecho}
% \def\documentsubtitle {Actividad 1 - Filosofia del Derecho}
% \def\documentsubject {Actividad Semana 2}
% ==============================================================================
\def\documenttitle {Reflexión y estudio de caso: interpretación jurídica}
\def\documentsubtitle {Actividad 5 - Filosofía del Derecho}
\def\documentsubject {Semana 7 - Interpretación jurídica}

\def\documentauthor {Martín Jonathan de la Cruz}
\def\coursename {Filosofía del Derecho}
\def\coursecode {FDE30}

\def\universityname {Universidad Abierta y a Distancia de México}
\def\universityfaculty {Licenciatura en Derecho}
\def\universitydepartment {Filosofía del Derecho}
\def\universitydepartmentimage {departamentos/UnADM}
\def\universitydepartmentimagecfg {height=1.57cm}
\def\universitylocation {Roma Norte, Ciudad de Mexico}

% INTEGRANTES, PROFESORES Y FECHAS
\def\authortable {
	\begin{tabular}{ll}
		\textbf{Alumno:}
		& \begin{tabular}[t]{l}
			 Martín Jonathan de la Cruz \\
		\end{tabular} \\
		Matrícula: & ES2611202040 \\

		\textit{Profesor:}
		& \begin{tabular}[t]{l}
			Reyna Patricia \\ González Rodríguez
		\end{tabular} \\
		& \\
		\multicolumn{2}{l}{Fecha de realización: \today} \\
		\multicolumn{2}{l}{\universitylocation}
	\end{tabular}
}

% IMPORTACION DEL TEMPLATE
\input{template}

% FORMATO SOLICITADO / USADO EN ACTIVIDADES CORREGIDAS
% La planeación solicita fuente Arial 12. En LaTeX clásico (pdfLaTeX) se usa Helvetica como sustituto:
\usepackage{helvet}
\renewcommand{\familydefault}{\sfdefault}
% Si desea usar Arial exactamente, compile con XeLaTeX o LuaLaTeX y descomente:
% \usepackage{fontspec}
% \setmainfont{Arial}
% Nota: verificar que el tamaño de letra sea 12pt en el documento final.

% PRODUCTOS VISUALES EN TIKZ
% Usar TikZ para construir productos solicitados cuando sea posible:
% cuadro comparativo, tabla, mapa conceptual, mapa mental, linea de tiempo,
% diagrama de flujo, mapa de cajas, redes conceptuales, esquemas, matrices,
% SQA, Venn, procesos, cronologias, arboles de decision y organizadores.
\usepackage{tikz}
\usetikzlibrary{
	arrows.meta,
	positioning,
	calc,
	matrix,
	fit,
	shapes.geometric,
	shadows.blur
}

% CITAS EN FORMATO AUTOR-ANIO, COMPATIBLE CON APA 7 EN EL CUERPO DEL TEXTO
% Usar:
% - Cita parentetica: \citep{clave}
% - Cita narrativa: \citet{clave}
% - Varias fuentes: \citep{clave1,clave2}
% - Parafrasis con pagina si conviene: \citep[p.~00]{clave}
% - Cita textual breve: "texto exacto" \citep[p.~00]{clave}
%
% Regla practica:
% - Si la consigna pide "con tus propias palabras", usar parafrasis con cita.
% - Si se toma una frase exacta, usar comillas y pagina.
% - En definiciones de glosario, citar inmediatamente despues de la definicion
%   base y desarrollar el analisis propio en las oraciones siguientes.
\setcitestyle{authoryear,open={(},close={)}}

% ==============================================================================
% AYUDAS DE REDACCION
% Quitar o comentar \pendiente cuando el documento final este terminado.
% ==============================================================================
\newcommand{\pendiente}[1]{\textcolor{red}{[PENDIENTE: #1]}}

% ==============================================================================
% MARCA DE AGUA EN PORTADA
% - Se aplica solo a la portada porque usa \AddToShipoutPictureBG*.
% - Para apagarla en alguna actividad: \def\coverwatermarkenabled {false}
% - Usar opacidad baja para que no compita con titulo, datos ni logotipo.
% ==============================================================================
\def\coverwatermarkenabled {true}
\def\coverwatermarkimage {img/departamentos/UnADM.pdf}
\def\coverwatermarkopacity {0.16}
\def\coverwatermarkwidth {0.72\paperwidth}
\def\coverwatermarkxshift {0cm}
\def\coverwatermarkyshift {-0.35cm}

\newcommand{\insertcoverwatermark}{%
	\ifthenelse{\equal{\coverwatermarkenabled}{true}}{%
		\AddToShipoutPictureBG*{%
			\begin{tikzpicture}[remember picture,overlay]
				\node[
					opacity=\coverwatermarkopacity,
					inner sep=0pt,
					xshift=\coverwatermarkxshift,
					yshift=\coverwatermarkyshift
				] at (current page.center) {%
					\includegraphics[width=\coverwatermarkwidth]{\coverwatermarkimage}%
				};
			\end{tikzpicture}%
		}%
	}{}%
}

\begin{document}

% PORTADA
\insertcoverwatermark
\templatePortrait

% CONFIGURACION DE PAGINA Y ENCABEZADOS
\templatePagecfg
\onehalfspacing

% ==============================================================================
% RESUMEN / ABSTRACT
% Debe ser breve, no copiar instrucciones.
% Formula recomendada:
% Tema + objetivo + enfoque/metodo + postura/resultado.
% ==============================================================================
\begin{abstractd}
	\textbf{Objetivo:} Analizar la interpretación jurídica y sus métodos (literal, histórica, teleológica y sistemática), incorporar hermenéutica y argumentación, y aplicar estos criterios al estudio de dos tesis de la Suprema Corte de Justicia de la Nación sobre el delito de violación.

	\textbf{Postura de análisis:} La interpretación jurídica no es un trámite técnico; es el punto donde el sistema decide qué derechos protege y qué límites impone al poder punitivo. Si la interpretación se ancla en estereotipos o formalismos, la norma pierde su función de garantía; si se sostiene con hermenéutica y argumentación, puede corregir sesgos y producir decisiones más justas.
\end{abstractd}

% TABLA DE CONTENIDOS - INDICE
\templateIndex

% CONFIGURACIONES FINALES
\templateFinalcfg

% ======================= INICIO DEL DOCUMENTO =======================

% ==============================================================================
% SECCION 1: INTRODUCCION
% Apertura recomendada:
% - Iniciar con una afirmacion fuerte.
% - Traducir la consigna a problema juridico.
% - Mostrar por que el tema importa.
% - Cerrar anticipando el enfoque personal.
%
% Evitar:
% - "En este trabajo hablare de..."
% - "Es muy importante..."
% - Definiciones de diccionario sin interpretacion.
% ==============================================================================
\section{Introducción}

Interpretar una norma no es solo leerla; es definir su sentido operativo dentro del sistema jurídico. En ese acto se decide qué derechos se protegen, qué conductas se consideran legítimas y qué límites se imponen al poder punitivo. Por eso la interpretación funciona como un criterio de control, no como un trámite mecánico.

En materia penal, la forma de interpretar conceptos como violencia o consentimiento determina si el sistema reconoce la vulnerabilidad de la víctima o reproduce estereotipos. Las dos tesis de la SCJN sobre el artículo 175 de los códigos penales de Jalisco y de la Ciudad de México muestran este giro interpretativo al ampliar el entendimiento de la violencia física y al descartar la incapacidad de resistencia como consentimiento \citep{scjnViolenciaFisica2022,scjnIncapacidadResistencia2019}.

La reflexión que sigue parte de los métodos de interpretación, la hermenéutica y la argumentación jurídica, y después aplica esos criterios al análisis de ambos precedentes. La posición de fondo es que interpretar exige justificar con razones y con efectos reales: lo que se decide en una sentencia termina definiendo la protección efectiva de los derechos.

\subsection*{Criterio de análisis}

Para evitar que el trabajo se convierta en un resumen de fuentes, el análisis se organiza en dos planos. Primero se fija una postura conceptual sobre la interpretación jurídica; después se revisa cada precedente a partir del texto normativo, los hechos, el problema jurídico, los criterios interpretativos, los argumentos y la consecuencia. Esa estructura ayuda a mantener una redacción personal: no se repite la tesis judicial, sino que se explica qué controla, qué corrige y qué riesgos interpretativos deja abiertos.

% ==============================================================================
% SECCION 2: DESARROLLO / ANALISIS
% Organizar segun el producto solicitado:
% - Mapa conceptual: concepto central, conceptos principales, conceptos secundarios,
%   relaciones y conclusion.
% - Cuadro comparativo: categorias, criterios, autores, fundamentos, principios,
%   caracteristicas, ejemplos, limites y valoracion critica.
% - Glosario: termino, definicion operativa, fuente, ejemplo juridico y comentario.
% - Estudio de caso: hechos, problema juridico, normas/principios, analisis,
%   postura y solucion.
%
% Regla derivada de Actividad 3:
% - Cuando la actividad pida glosario y analisis de caso, no basta con enumerar
%   conceptos. Ubicar primero la funcion de los conceptos y despues explicar cada
%   uno con definicion base + cita + comentario propio.
% - En la tabla del caso, no repetir la definicion del glosario: aplicar el
%   concepto a hechos, norma, criterio judicial, consecuencia o tension moral.
%
% Regla de parrafo:
% Entrada: idea/concepto.
% Proceso: como funciona.
% Salida: por que importa o que consecuencia produce.
% ==============================================================================
\section{Reflexión}

La interpretación jurídica es el proceso mediante el cual se atribuye sentido a las normas para resolver problemas concretos; no es un acto neutral, sino una decisión que conecta texto, contexto y consecuencias. Por eso la interpretación opera como un filtro de validez práctica: define cómo se protege un derecho y qué conductas se sancionan \citep{hernandezManriquezHermeneutica2019}.

La importancia de interpretar correctamente una norma dentro del sistema jurídico mexicano radica en que el texto legal, por sí solo, no resuelve todos los casos. Hernández Manríquez sostiene que la hermenéutica jurídica permite comprender el Derecho desde su contexto, sus fines y sus efectos, no solo desde la literalidad de las palabras \citep{hernandezManriquezHermeneutica2019}. Retomo esa idea porque, en el contexto mexicano, muchas controversias jurídicas aparecen precisamente cuando una lectura gramatical resulta insuficiente para proteger derechos o para controlar abusos de poder. Desde mi perspectiva, interpretar es decidir qué peso tendrá la norma en la vida real: si se aplica de forma rígida, puede producir injusticia; si se interpreta con método, puede convertirse en una herramienta de garantía.

\subsection{Métodos de interpretación}

\begin{itemize}
	\item \textbf{Interpretación literal.} Parte del significado gramatical de la norma. Su ventaja es la seguridad jurídica, pero su límite es el formalismo cuando el texto se vuelve insuficiente para captar el daño real \citep{scjnMemoriaArgumentacion2008}.
	\item \textbf{Interpretación histórica.} Busca la intención del legislador y el contexto en que nació la norma. Ayuda a fijar alcance, pero puede congelar el sentido y dejar fuera transformaciones sociales \citep{hernandezManriquezHermeneutica2019}.
	\item \textbf{Interpretación teleológica.} Atiende al fin de la norma. Permite recuperar su función protectora, aunque requiere argumentar con cuidado para evitar lecturas arbitrarias \citep{scjnMemoriaArgumentacion2008}.
	\item \textbf{Interpretación sistemática.} Integra la norma dentro del sistema jurídico y sus principios. Da coherencia y evita contradicciones, especialmente cuando se conecta con derechos humanos \citep{scjnIncapacidadResistencia2019,scjnViolenciaFisica2022}. En particular, la interpretación sistemática debe articular mandatos constitucionales y normas de protección de víctimas (p. ej. \citep{cpeum2026,LeyGeneralVictimas}), garantizando que la lectura de la norma responda a obligaciones constitucionales y convencionales.
\end{itemize}

La hermenéutica jurídica amplía la mirada al reconocer que toda interpretación tiene un horizonte de sentido: no basta traducir palabras, hay que comprender prácticas, estructuras de poder y efectos sobre las personas \citep{hernandezManriquezHermeneutica2019}. La argumentación jurídica, por su parte, exige justificar por qué un método es preferible a otro en un caso concreto; sin esa justificación, la interpretación pierde legitimidad y se vuelve discrecional \citep{scjnMemoriaArgumentacion2008}. Desde esta perspectiva, interpretar bien no es elegir el método más cómodo, sino el que resuelve el problema con mayor coherencia, control y protección de derechos.

En mi lectura, ninguno de estos métodos funciona de manera aislada. La interpretación literal ofrece un punto de partida necesario, pero no suficiente; la histórica ayuda a entender por qué nació la norma; la teleológica permite recuperar su finalidad protectora; y la sistemática evita que una disposición se aplique al margen de la Constitución y de los derechos humanos. Por eso considero que, en el sistema jurídico mexicano, la interpretación sistemática y la teleológica tienen especial relevancia cuando están respaldadas por una base literal suficiente y por una argumentación rigurosa. Su valor no consiste en sustituir el texto, sino en impedir que el texto se convierta en una excusa para decisiones injustas.

En consecuencia, la importancia de la interpretación jurídica no radica solo en aclarar qué dice una norma, sino en definir cómo operará frente a conflictos reales. Una mala interpretación puede normalizar estereotipos, vaciar derechos o volver insuficiente la protección de la víctima; una interpretación justificada, en cambio, permite conectar la ley con la dignidad humana, la igualdad y el acceso efectivo a la justicia. Esa es la razón por la que sostengo que interpretar, dentro del sistema jurídico mexicano, implica una responsabilidad argumentativa y no una simple operación gramatical.

% ==============================================================================
% BLOQUE OPCIONAL: GLOSARIO Y ANALISIS DE CASO
% Usar cuando la planeacion pida "Glosario", "conceptos juridicos fundamentales",
% "analisis de caso" o una tabla concepto/aplicacion practica.
%
% Reglas derivadas de la retroalimentacion:
% - Minimo de conceptos: respetar el numero de la planeacion.
% - Orden alfabetico si se solicita.
% - Cada concepto: definicion propia + cita APA + funcion juridica.
% - No usar citas textuales en todos los conceptos si se piden palabras propias.
% - Cita textual solo si es indispensable; incluir pagina.
% - Antes del glosario, se puede incluir una ubicacion conceptual breve, no una
%   explicacion larga de la consigna.
% - En el caso: identificar al menos los conceptos pedidos y explicar como operan.
% - Conclusion: actos juridicos identificados + relacion con moral/derecho +
%   postura personal.
% ==============================================================================
%
% \section{Glosario de conceptos juridicos fundamentales}
%
% \subsection*{Ubicacion conceptual previa}
% \pendiente{Clasificar brevemente los conceptos por funcion: sistema juridico,
% norma, sujetos, relacion juridica, hechos/actos, moral/derecho. No copiar la
% instruccion de la planeacion.}
%
% \begin{description}
% \item[\textbf{Concepto.}]
% Definicion base en palabras propias \citep{claveFuente}. Funcion del concepto
% dentro del sistema juridico. Ejemplo breve o consecuencia si se interpreta mal.
%
% \item[\textbf{Otro concepto.}]
% Definicion base en palabras propias \citep[p.~00]{claveFuente}. Comentario
% propio: que problema resuelve, que limite tiene y como se aplica.
% \end{description}
%
% \section{Analisis de caso}
%
% \begingroup
% \small
% \setlength{\tabcolsep}{5pt}
% \renewcommand{\arraystretch}{1.35}
% \begin{longtable}{@{}p{0.24\linewidth}p{0.70\linewidth}@{}}
% \caption{Conceptos juridicos aplicados al caso}\label{tab:analisis-caso}\\
% \toprule
% \textbf{Concepto} & \textbf{Aplicacion practica en el caso} \\
% \midrule
% \endfirsthead
% \toprule
% \textbf{Concepto} & \textbf{Aplicacion practica en el caso} \\
% \midrule
% \endhead
% \bottomrule
% \endlastfoot
%
% \textbf{Concepto juridico} &
% Aplicar el concepto al hecho, norma, autoridad, derecho, deber, sancion,
% validez, eficacia o tension moral del caso. No repetir la definicion. \\
% \midrule
%
% \end{longtable}
% \endgroup

% ==============================================================================
% BLOQUE OPCIONAL: CUADRO COMPARATIVO
% Usar cuando la actividad pida cuadro comparativo.
% Comentarios de rubrica:
% - Categorias claras y jerarquizadas.
% - No solo datos: incluir funcion, limites, aplicacion y valoracion critica.
% - Al final del cuadro debe venir conclusion/postura personal si la planeacion lo pide.
% ==============================================================================
%
% \clearpage
% \pagestyle{empty}
% \begin{landscape}
% \begingroup
% \scriptsize
% \setlength{\tabcolsep}{4pt}
% \renewcommand{\arraystretch}{1.35}
% \begin{longtable}{@{}p{0.18\linewidth}p{0.39\linewidth}p{0.39\linewidth}@{}}
% \caption{Cuadro comparativo entre [tema A] y [tema B]}\label{tab:cuadro-comparativo}\\
% \toprule
% \textbf{Criterio} & \textbf{Tema A} & \textbf{Tema B} \\
% \midrule
% \endfirsthead
% \toprule
% \textbf{Criterio} & \textbf{Tema A} & \textbf{Tema B} \\
% \midrule
% \endhead
% \midrule
% \multicolumn{3}{r}{Continua en la siguiente pagina} \\
% \endfoot
% \bottomrule
% \endlastfoot
%
% Concepto &
% [Definicion con funcion, no solo descripcion. Cita APA.] &
% [Definicion con funcion, no solo descripcion. Cita APA.] \\
% \midrule
%
% Fundamento &
% [Base teorica, juridica o filosofica.] &
% [Base teorica, juridica o filosofica.] \\
% \midrule
%
% Principios &
% [Principios ordenados por jerarquia.] &
% [Principios ordenados por jerarquia.] \\
% \midrule
%
% Aplicacion &
% [Ejemplo real o normativo.] &
% [Ejemplo real o normativo.] \\
% \midrule
%
% Valoracion critica &
% [Fortaleza, limite y costo si se aplica mal.] &
% [Fortaleza, limite y costo si se aplica mal.] \\
%
% \end{longtable}
% \endgroup
% \end{landscape}
% \pagestyle{fancy}

% ==============================================================================
% BLOQUE OPCIONAL: MAPA CONCEPTUAL O FIGURA
% Usar si la actividad pide mapa conceptual, esquema o evidencia visual.
% Si se usa herramienta externa, insertar captura y, si procede, enlace.
%
% Reglas por retroalimentacion:
% - Si la planeacion pide mapa conceptual, el producto principal debe ser mapa,
%   no un texto expositivo disfrazado.
% - Incluir concepto central claramente identificado.
% - Incluir conceptos principales y secundarios.
% - Usar palabras de enlace en las flechas o conexiones.
% - Mostrar jerarquia logica: no repetir conceptos en el mismo nivel sin relacion.
% - Agregar conclusion breve solo si la planeacion lo pide o ayuda a cerrar.
% ==============================================================================
%
% \clearpage
% \begin{figure}[H]
% 	\centering
% 	\includegraphics[width=0.95\linewidth]{img/nombre-del-mapa.png}
% 	\caption{Mapa conceptual de [tema]}
% 	\label{fig:mapa-conceptual}
% \end{figure}

% ==============================================================================
% SECCION 3: POSTURA / ANALISIS PERSONAL
% Esta seccion evita la observacion de "informacion general" o "no se detecta
% autenticidad". Debe contestar:
% - Que posicion se sostiene.
% - Por que esa posicion es razonable.
% - Que evidencia/fuente la respalda.
% - Como se aplica al sistema juridico mexicano.
% ==============================================================================
\section{Análisis de caso}

\subsection{Criterio 1: violencia física en el delito de violación (Jalisco)}

\textbf{Texto normativo de partida.} El artículo 175 del Código Penal para el Estado de Jalisco sanciona a quien, por medio de violencia física o moral, tenga cópula con una persona mayor de edad. Ese punto de partida es decisivo porque toda la controversia gira en torno al alcance que debe darse a la expresión \enquote{violencia física}: si debe leerse en sentido estrecho, como fuerza corporal directa, o en sentido funcional, como imposición material sobre una persona colocada en estado de vulnerabilidad \citep{scjnViolenciaFisica2022}.

\textbf{Hechos.} En el caso analizado, la imputación por violación se sustentó en que la víctima perdió la consciencia tras ingerir alcohol y el imputado, aprovechándose de esa situación, la trasladó a otro lugar y le impuso la cópula. El juez de control dictó un auto de no vinculación a proceso al considerar que no se acreditaba violencia física; sin embargo, la instancia revisora revocó esa decisión. El criterio judicial sostuvo que la violencia física no exige necesariamente fuerza visible ni resistencia explícita de la víctima \citep{scjnViolenciaFisica2022}.

\textbf{Normas jurídicas aplicadas.} La norma penal directamente aplicable es el artículo 175 del Código Penal para el Estado de Jalisco, porque define el tipo de violación y exige violencia física o moral como medio de comisión. Además, la interpretación del caso debe leerse en armonía con el artículo 1 constitucional y con la Ley General de Víctimas, ya que la protección de la integridad y de la dignidad de la víctima impide interpretar el tipo penal desde estereotipos sobre resistencia física o conducta esperada \citep{cpeum2026,LeyGeneralVictimas}.

\textbf{Problema jurídico.} Determinar si la violencia física debe limitarse a fuerza corporal directa y resistencia visible, o si también comprende el aprovechamiento de una condición de vulnerabilidad que impide a la víctima reaccionar. El punto crítico es que una lectura estrecha puede dejar fuera conductas que afectan la libertad sexual aunque no presenten la imagen tradicional de agresión.

\textbf{Interpretación literal.} La lectura gramatical podría entender la violencia física como fuerza corporal directa y resistencia de la víctima; bajo esta lectura, el caso quedaría fuera si no se acredita oposición física.

\textbf{Interpretación histórica.} Las concepciones tradicionales del delito exigían fuerza y resistencia como signos visibles de violencia; este antecedente explica por qué el juez de control adoptó un criterio restrictivo.

\textbf{Interpretación teleológica.} El fin de la norma es proteger la libertad sexual y la integridad personal; si el agresor se aprovecha de un estado de vulnerabilidad, se actualiza la violencia física como medio de imposición, aun sin golpes o resistencia.

\textbf{Interpretación sistemática.} La norma debe leerse de forma coherente con el sistema de derechos humanos y con la protección de víctimas; por ello, la violencia física se entiende de manera amplia cuando existe control y aprovechamiento de la vulnerabilidad.

\textbf{Postura y consecuencia.} Desde mi lectura, el criterio es correcto porque desplaza el análisis desde la reacción de la víctima hacia la conducta del agresor. Esa diferencia no es menor: si el sistema exige resistencia visible, termina castigando a la víctima por no haber reaccionado; si observa el control ejercido sobre ella, la interpretación recupera la finalidad protectora de la norma. El límite está en exigir prueba suficiente sobre la condición de vulnerabilidad y sobre el aprovechamiento de esa condición.

\textbf{Argumentos a favor.}
\begin{enumerate}
	\item \textbf{Corrige una lectura estereotipada del tipo penal.} El principal acierto del criterio es que rompe con la idea de que la violencia solo existe cuando hay golpes visibles o resistencia física. Esa exigencia tradicional no solo es limitada, sino que puede distorsionar la función de la norma, porque desplaza la atención desde el acto del agresor hacia el comportamiento esperado de la víctima. La tesis corrige ese sesgo y reubica el problema donde corresponde: en la imposición sexual ejercida sobre una persona en situación de vulnerabilidad.
	\item \textbf{Fortalece la tutela de la libertad sexual.} La interpretación amplia de violencia física resulta razonable porque protege de mejor manera el bien jurídico tutelado. Si la libertad sexual puede lesionarse mediante el aprovechamiento de una incapacidad momentánea, entonces una lectura demasiado estrecha del tipo penal dejaría sin respuesta conductas materialmente gravosas. En ese sentido, el criterio no amplía por capricho, sino para conservar la eficacia protectora de la norma.
	\item \textbf{Mejora el acceso a la justicia.} Exigir resistencia corporal como prueba central suele invisibilizar agresiones cometidas en contextos de inconsciencia o debilidad extrema. El criterio reduce ese obstáculo probatorio y permite una valoración más adecuada de los hechos. Esto es importante porque la justicia penal no debe depender de estereotipos sobre cómo debe reaccionar una víctima, sino de la prueba sobre el control ejercido por el agresor y sobre la afectación real de la libertad sexual.
\end{enumerate}

\textbf{Argumentos en contra.}
\begin{enumerate}
	\item \textbf{Riesgo de expansión interpretativa excesiva.} Una objeción posible es que el concepto de violencia física se vuelva demasiado amplio si no se delimitan con precisión las condiciones de vulnerabilidad relevantes. Si toda situación de desventaja se asimila automáticamente a violencia física, el tipo penal puede perder claridad y volverse menos previsible.
	\item \textbf{Tensión con el principio de legalidad.} También puede sostenerse que una redefinición judicial demasiado intensa del concepto de violencia física corre el riesgo de aproximarse a una labor legislativa. En materia penal, esa objeción es seria porque el principio de taxatividad exige que los elementos típicos sean suficientemente claros y no dependan de reconstrucciones abiertas sin control argumentativo.
	\item \textbf{Necesidad de estándares probatorios consistentes.} El criterio solo será legítimo si se acompaña de parámetros claros para demostrar la vulnerabilidad y el aprovechamiento de esa condición. De lo contrario, la amplitud interpretativa podría generar incertidumbre jurídica y decisiones desiguales. Por eso, la ventaja protectora de la tesis exige un complemento indispensable: rigor probatorio y justificación judicial suficiente.
\end{enumerate}

\subsection{Criterio 2: incapacidad de resistencia y consentimiento en la violación equiparada (CDMX)}

\textbf{Texto normativo de partida.} El artículo 175 del Código Penal del Distrito Federal, aplicable para la Ciudad de México, contempla la violación equiparada cuando la persona ofendida no puede resistir o no comprende el significado del hecho. El centro del problema interpretativo no está en repetir la fórmula legal, sino en establecer qué debe entenderse por incapacidad relevante para negar valor jurídico al consentimiento \citep{scjnIncapacidadResistencia2019}.

\textbf{Hechos.} La tesis analiza un caso de violación equiparada en el que el problema no consistió en probar violencia física o moral, sino en determinar si la condición de la víctima impedía una decisión sexual verdaderamente voluntaria. El criterio sostiene que cuando la persona no puede resistir o no comprende lo que ocurre, esa situación no puede leerse como consentimiento, porque consentir exige capacidad real de decisión \citep{scjnIncapacidadResistencia2019}.

\textbf{Normas jurídicas aplicadas.} La norma penal directamente aplicable es el artículo 175 del Código Penal del Distrito Federal, porque regula la violación equiparada cuando la persona no puede resistir o no comprende el significado del hecho. A ello se suman el artículo 1 constitucional y la Ley General de Víctimas, ya que la interpretación del consentimiento debe ser compatible con la dignidad humana, la igualdad y la protección reforzada de personas en situación de vulnerabilidad \citep{cpeum2026,LeyGeneralVictimas}.

\textbf{Problema jurídico.} Definir si la ausencia de resistencia puede interpretarse como consentimiento, o si el consentimiento requiere capacidad real de comprender, decidir y expresar una voluntad libre. El punto más difícil está en determinar qué debe entenderse por incapacidad suficiente para consentir, porque no siempre coincide con una pérdida total de consciencia. Una persona puede encontrarse afectada física o mentalmente sin quedar completamente inconsciente, y esa zona intermedia obliga a distinguir entre simple alteración, vulnerabilidad real e imposibilidad efectiva de autodeterminación. La consecuencia jurídica es directa: si se confunde pasividad con consentimiento, se reduce la protección penal; pero si se presume incapacidad sin elementos objetivos, también se corre el riesgo de vaciar indebidamente la autonomía personal.

\textbf{Interpretación literal.} El texto legal distingue entre violación básica y equiparada; en esta última no se exige violencia física o moral, sino incapacidad de resistencia o comprensión.

\textbf{Interpretación histórica.} El legislador incorporó la figura para proteger a personas en condiciones de vulnerabilidad que no pueden oponerse o comprender, ampliando la tutela penal frente a agresiones sexuales.

\textbf{Interpretación teleológica.} El propósito es proteger la libertad y autonomía sexual, evitando que la pasividad o la incapacidad se interpreten como aceptación válida. Desde esta finalidad, el consentimiento no puede reducirse a la mera ausencia de negativa visible; exige una aptitud real para comprender el acto y decidir sobre él.

\textbf{Interpretación sistemática.} La norma se integra con principios de dignidad, igualdad y protección reforzada de personas vulnerables; por ello, el consentimiento solo es válido cuando existe capacidad real de decisión. Esta lectura obliga a valorar conjuntamente el estado cognitivo, la posibilidad de comprender lo que ocurre, la facultad de expresar voluntad y el contexto concreto en que sucedieron los hechos.

\textbf{Postura y consecuencia.} La incapacidad para consentir en delitos sexuales constituye uno de los problemas interpretativos más complejos del Derecho penal, porque no siempre coincide con una pérdida total de conciencia. La dificultad central consiste en determinar cuándo una persona carece realmente de aptitud para decidir libremente sobre su conducta sexual.

No basta afirmar que alguien "estaba mal" o "estaba ebrio"; jurídicamente debe acreditarse si esa condición anuló su posibilidad de comprender el acto, valorar sus consecuencias y expresar una voluntad libre. En este punto, la interpretación no puede reducirse a un criterio exclusivamente físico (por ejemplo, fuerza o inmovilidad), ni tampoco a una etiqueta genérica de "alteración mental". La evaluación debe considerar, de manera conjunta, el estado cognitivo, la capacidad volitiva y el contexto del hecho. Desde esta perspectiva, la ebriedad no opera automáticamente como consentimiento ni como ausencia de consentimiento en todos los casos: su relevancia jurídica depende del grado de afectación real de las facultades mentales y de la posibilidad efectiva de autodeterminación. Por ello, el estándar probatorio debe orientarse a verificar si existió una disminución relevante o anulación de la capacidad para decidir, y no únicamente la presencia de alcohol u otra condición clínica.

Además, el análisis debe distinguir entre discapacidad, trastorno mental, intoxicación transitoria y estados de inconsciencia, porque cada supuesto plantea niveles distintos de vulnerabilidad y de protección jurídica. Confundir esas categorías puede producir decisiones injustas: o bien se invalida indebidamente la voluntad de personas capaces, o bien se exige una resistencia imposible a personas que no podían consentir de forma válida.

La postura que sostengo es que la interpretación constitucional y penal sobre consentimiento debe centrarse en la autonomía real de la persona y en la protección reforzada frente a situaciones de vulnerabilidad. En consecuencia, la "incapacidad para consentir" debe entenderse como la imposibilidad material y/o psíquica de formar, expresar o sostener una voluntad libre, demostrada con criterios objetivos, periciales y contextuales, no con prejuicios ni con estereotipos sobre cómo "debería" comportarse una víctima.

\textbf{Argumentos a favor.}
\begin{enumerate}
	\item \textbf{Diferencia consentimiento de simple pasividad.} El mayor mérito del criterio es que aclara que la ausencia de oposición visible no equivale a una aceptación válida. Esa precisión es central en materia penal sexual, porque evita que la carga del análisis recaiga en demostrar resistencia física en lugar de acreditar la existencia de una voluntad libre e informada.
	\item \textbf{Protege mejor a personas vulnerables.} La tesis reconoce que hay contextos en los que la persona puede encontrarse incapacitada para comprender o reaccionar sin estar totalmente inconsciente. Este matiz es importante porque la vulnerabilidad real no siempre adopta la forma de inconsciencia absoluta. Al admitir esa complejidad, la interpretación responde mejor a la experiencia concreta de los hechos y no solo a una imagen simplificada del consentimiento.
	\item \textbf{Se alinea con un enfoque de derechos humanos.} Exigir capacidad real de decisión es compatible con la dignidad humana y con la protección reforzada de la autonomía sexual. La tesis fortalece ese estándar al afirmar que el consentimiento debe ser una manifestación libre y no una presunción derivada del silencio, de la inmovilidad o de la confusión de la víctima.
\end{enumerate}

\textbf{Argumentos en contra.}
\begin{enumerate}
	\item \textbf{Riesgo de paternalismo interpretativo.} Si la tesis se aplica sin distinguir cuidadosamente entre discapacidad, intoxicación transitoria, trastorno mental o simple alteración emocional, podría terminar sustituyendo indebidamente la autonomía de personas que sí conservaban capacidad de decidir. La protección penal no debe convertirse en negación automática de agencia personal.
	\item \textbf{Ambigüedad en la frontera de la incapacidad.} Otra objeción relevante es que la línea entre consentimiento válido e incapacidad suficiente puede volverse difusa cuando no existen parámetros claros. Ese problema se vuelve más visible en casos de ebriedad, somnolencia, medicación o afectaciones psicológicas parciales, donde la valoración exige mayor precisión técnica.
	\item \textbf{Necesidad de control probatorio reforzado.} Un estándar amplio de incapacidad puede generar incertidumbre jurídica si no se acompaña de prueba objetiva y de una motivación judicial rigurosa. Por ello, la tesis es correcta en su finalidad protectora, pero necesita un uso prudente: el juez debe demostrar por qué, en el caso concreto, la persona no podía comprender, decidir o sostener una voluntad libre.
\end{enumerate}

% ==============================================================================
% SECCION 4: CONCLUSION
% Debe cerrar con fuerza:
% - Sintetizar el hallazgo.
% - Reafirmar postura.
% - Indicar aplicacion real.
% - Evitar repetir la introduccion.
%
% Si la actividad pregunta "cual corriente tiene mayor relevancia y por que",
% responder de forma directa, no dejarlo implicito.
% ==============================================================================

% ========================= CUMPLIMIENTO DE LA RÚBRICA ==========================
% A continuación se indica de forma explícita cómo el trabajo satisface cada
% criterio señalado en la planificación y en la retroalimentación editorial.
% Esto facilita la revisión y reduce las posibilidades de descuento por
% incumplimiento formal o de contenido.

\section*{Cumplimiento de la rúbrica}

\begin{itemize}
	\item \textbf{Portada:} Generada por la plantilla institucional (ver portada).
	\item \textbf{Introducción:} Existe y plantea el problema central de la actividad
	\quad (sección \textit{Introducción}).
	\item \textbf{Reflexión:} Desarrollo teórico y postura personal sobre la
	\quad interpretación jurídica (sección \textit{Reflexión}).
	\item \textbf{Análisis de caso:} Ambos precedentes de la SCJN son analizados por
	\quad separado y con la misma estructura exigida en la planeación: hechos,
	\quad norma aplicable, problema jurídico, interpretación literal/histórica/
	\quad teleológica/sistemática, argumentos a favor y en contra (subsecciones
	\quad "Criterio 1" y "Criterio 2").
	\item \textbf{Conclusiones:} Dos párrafos de cierre con postura final y
	\quad enseñanza práctica (sección \textit{Conclusiones}).
	\item \textbf{Referencias:} Bibliografía en archivo \texttt{filosofia-del-derecho.bib}
	\quad y citada con natbib (\texttt{\setcitestyle{authoryear}}) para cumplir
	\quad formato autor-año compatible con APA.
    \item \textbf{Formato y presentación:} Se usó la plantilla institucional
    \quad (Helvética/Arial equivalencia), interlineado 1.5, títulos en negritas y
	\quad texto justificado según las reglas de la planeación.
	\item \textbf{Producto y extensión:} El trabajo está organizado según la
	\quad estructura solicitada y evita copiar instrucciones de la planeación en
	\quad el cuerpo del documento.
	\item \textbf{Originalidad y postura propia:} Se adicionó y reforzó una sección
	\quad de postura argumentada y se transformaron observaciones de retroalimentación
	\quad en redacción académica (evita tono de resumen automático).
\end{itemize}

% Checklist operativo (marcar en la revisión final antes de entregar):
\begin{itemize}
	\item[\checkmark] Portada institucional correcta
	\item[\checkmark] Introducción clara con problema central
	\item[\checkmark] Reflexión teórica con postura propia
	\item[\checkmark] Análisis de casos con estructura pedida
	\item[\checkmark] Tres argumentos a favor y tres en contra por criterio
	\item[\checkmark] Conclusiones (2 párrafos) y cierre académico
	\item[\checkmark] Bibliografía (.bib) coincide con las citas
	\item[\checkmark] Formato (fuente, interlineado, negritas)
\end{itemize}

% Observaciones finales sobre requisitos formales:
% - Verificar que el PDF final respete la nomenclatura solicitada: 
%   FDE30_ApellidoPaterno_Nombre_Act5.pdf
% - Confirmar que la compilación genere la bibliografía completa (ejecutar
%   latexmk o el script provisto en \texttt{scripts/}).

\section{Conclusiones}

La actividad permitió confirmar que la interpretación jurídica no es una operación secundaria del Derecho, sino el mecanismo que define su eficacia real. La reflexión conceptual mostró que interpretar una norma exige articular literalidad, contexto, finalidad y sistema; por eso, la hermenéutica y la argumentación resultan indispensables cuando el texto legal, por sí solo, no basta para resolver conflictos complejos. En los dos criterios analizados, la SCJN desplaza el foco desde la resistencia física o la pasividad aparente hacia la protección efectiva de la autonomía sexual, corrigiendo estereotipos que durante mucho tiempo limitaron el reconocimiento jurídico de la violencia \citep{scjnViolenciaFisica2022,scjnIncapacidadResistencia2019}.

La postura final que sostengo es que, en el sistema jurídico mexicano, la interpretación sistemática y la teleológica adquieren especial importancia cuando está en juego la dignidad de la persona y la protección de víctimas, siempre que permanezcan controladas por el texto legal y por una argumentación rigurosa. En términos prácticos, esto significa que interpretar bien no equivale a ampliar arbitrariamente la ley, sino a justificar por qué una lectura protege mejor los derechos sin vaciar el principio de legalidad. Ese equilibrio es el punto central del análisis realizado y, a mi juicio, la principal enseñanza jurídica de esta actividad.\footnote{En la elaboración de este trabajo se utilizaron herramientas de inteligencia artificial como apoyo para la organización y revisión del texto, conforme a los Lineamientos para el uso ético y responsable de la Inteligencia Artificial en el ámbito académico de la UnADM. El análisis, la selección de fuentes y la postura sostenida son de mi autoría.}

% ==============================================================================
% REFERENCIAS
% Reglas:
% - Toda fuente citada en el texto debe estar en el .bib.
% - Toda referencia en el .bib debe estar citada en el texto final; quitar
%   fuentes que no se usaron.
% - Toda cita debe tener clave correcta y coincidir con autor/anio visible.
% - Usar fuentes de la planeacion y, si se complementa, que sean confiables.
% - En el cuerpo: \citep{clave}; no usar solo URLs sueltas.
% - Si se trabaja con PDF local, registrar datos bibliograficos completos:
%   autor, anio, titulo, editorial/institucion y URL/DOI si existe.
% - Si la fuente no tiene fecha, usar s. f. con cuidado y preferir confirmar los
%   datos en la portada o ficha bibliografica del documento.
% ==============================================================================

\bibliography{filosofia-del-derecho-clean}

% FIN DEL DOCUMENTO
\end{document}
